\setcounter{section}{1}


  \zwischenueberschrift{Übungsaufgaben}


\inputaufgabe
{}
{
In der linearen Algebra geht es oft um
\definitionsverweis {lineare Gleichungssysteme}{}{.}
Welche nichtlinearen Gleichungen bzw. Gleichungssysteme treten in der linearen Algebra auf?

}
{} {}


\inputaufgabe
{}
{
Skizziere im $\R^2$ die Lösungsmengen der folgenden Gleichungen.
\aufzaehlungzweireihe {\itemfuenf {
\mathl{x^2-y^2 -1 = 0}{,}
}{
\mathl{x^2+xy+y^2 = 0}{,}
}{
\mathl{x^2+y^2 +1 = 0}{,}
}{
\mathl{x^2+y^2 = 0}{,}
}{
\mathl{x^2+y^3 = 0}{,}
} } {\itemfuenf  {
\mathl{x^3-y^5 = 0}{,}
}{
\mathl{x^2-x^3 = 0}{,}
}{
\mathl{x^3+y^3 = 1}{,}
}{
\mathl{x^4+y^4 = 1}{,}
}{
\mathl{-5+3x+4x^2+x^3-y^2 = 1}{.}
} }

}
{} {}


\inputaufgabe
{}
{
Berechne den Durchschnitt der Kurven aus
Aufgabe 1.2
mit den folgenden Geraden.
\aufzaehlungsieben{
\mathl{x=0}{,}
}{
\mathl{y=0}{,}
}{
\mathl{x=1}{,}
}{
\mathl{y=-2}{,}
}{
\mathl{x= y}{,}
}{
\mathl{x=-y}{,}
}{
\mathl{2x-3y+4=0}{.}
}

}
{} {}


\inputaufgabegibtloesung
{}
{
\aufzaehlungzwei {Finde eine ganzzahlige Lösung

\mavergleichskette
{\vergleichskette
{(x,y)
}
{ \in }{ \Z \times \Z
}
{  }{
}
{    }{
}
{   }{
}
}
{}{}{}
für die Gleichung

\mavergleichskettedisp
{\vergleichskette
{ x^2-y^3+2
}
{ =} { 0
}
{ } {
}
{ } {
}
{ } {
}
}
{}{}{.}
} {Zeige, dass

\mathdisp {\left(  { \frac{   383 }{  1000 } }   , \,   { \frac{   129 }{  100 } }                   \right)} {  }

eine Lösung für die Gleichung

\mavergleichskettedisp
{\vergleichskette
{ x^2-y^3+2
}
{ =} { 0
}
{ } {
}
{ } {
}
{ } {
}
}
{}{}{}
ist.
}

}
{} {}


\inputaufgabegibtloesung
{}
{
Finde auf der
\definitionsverweis {ebenen algebraischen Kurve}{}{}

\mavergleichskettedisp
{\vergleichskette
{V(X^3-Y^3+4X^2-2XY+Y+3)
}
{ \subset} { {\mathbb C}^2
}
{ } {
}
{ } {
}
{ } {
}
}
{}{}{}
einen Punkt.

}
{} {}


\inputaufgabe
{}
{
Es sei $K$ ein
\definitionsverweis {Körper}{}{.}
Das
\definitionsverweis {Bild}{}{}
der durch
\maabbeledisp {} { K } { K^2
} { t } { \left( t^2  , \,  t^3                   \right)
} {,}
definierten Kurve heißt \stichwort {Neilsche Parabel} {.} Zeige, dass ein Punkt

\mavergleichskette
{\vergleichskette
{ (x,y)
}
{ \in }{ K^2
}
{  }{
}
{    }{
}
{   }{
}
}
{}{}{}
genau dann zu diesem Bild gehört, wenn er die Gleichung

\mavergleichskette
{\vergleichskette
{ x^3
}
{ = }{ y^2
}
{  }{
}
{    }{
}
{   }{
}
}
{}{}{}
erfüllt.

}
{} {}


\inputaufgabe
{}
{
Es sei

\mavergleichskette
{\vergleichskette
{ C
}
{ \subseteq }{ \R^2
}
{  }{
}
{    }{
}
{   }{
}
}
{}{}{}
das
\definitionsverweis {Bild}{}{}
unter der polynomialen Abbildung
\maabbeledisp {} { \R } { \R^2
} { t } { \left( t^2-t+2  , \,  t^2-3                   \right)
} {.}
Bestimme ein Polynom

\mavergleichskette
{\vergleichskette
{ F
}
{ \neq }{ 0
}
{  }{
}
{    }{
}
{   }{
}
}
{}{}{}
in zwei Variablen derart, dass $C$ auf dem Nullstellengebilde zu $F$ liegt.

}
{} {}


\inputaufgabe
{}
{
Es sei

\mavergleichskette
{\vergleichskette
{ C
}
{ \subseteq }{ \R^2
}
{  }{
}
{    }{
}
{   }{
}
}
{}{}{}
das
\definitionsverweis {Bild}{}{}
unter der polynomialen Abbildung
\maabbeledisp {} { \R } { \R^2
} { t } { \left( t^3-1  , \,  t^2-1                   \right)
} {.}
Bestimme ein Polynom

\mavergleichskette
{\vergleichskette
{ F
}
{ \neq }{ 0
}
{  }{
}
{    }{
}
{   }{
}
}
{}{}{}
in zwei Variablen derart, dass $C$ auf dem Nullstellengebilde zu $F$ liegt.

}
{} {}


\inputaufgabe
{}
{
Wir betrachten die Kurve
\maabbeledisp {} {\R} {\R^2
} { t } {(t^2-1,t^3-t)
} {.}
\aufzaehlungdreiabc{Zeige, dass die
\definitionsverweis {Bildpunkte}{}{}

\mathl{(x,y)}{} der Kurve die Gleichung

\mavergleichskettedisp
{\vergleichskette
{ y^2
}
{ =} { x^2+x^3
}
{ } {
}
{ } {
}
{ } {
}
}
{}{}{}
erfüllen.
}{Zeige, dass jeder Punkt

\mavergleichskette
{\vergleichskette
{ (x,y)
}
{ \in }{ \R^2
}
{  }{
}
{    }{
}
{   }{
}
}
{}{}{}
mit

\mavergleichskette
{\vergleichskette
{ y^2
}
{ = }{ x^2+x^3
}
{  }{
}
{    }{
}
{   }{
}
}
{}{}{}
zum Bild der Kurve gehört.
}{Zeige, dass es genau zwei Punkte
\mathkor {} {t_1} {und} {t_2} {}
mit identischem Bildpunkt gibt, und dass ansonsten die Abbildung injektiv ist.
}

}
{} {}


\inputaufgabe
{}
{
Diskutiere den Zusammenhang zwischen
\definitionsverweis {ebenen algebraischen Kurven}{}{}
und
dem Satz über implizite Funktionen.

}
{} {}


\inputaufgabe
{}
{
Es sei

\mavergleichskette
{\vergleichskette
{ K
}
{ = }{ \Z/( 7 )
}
{  }{
}
{    }{
}
{   }{
}
}
{}{}{.}
Bestimme alle Punkte in

\mavergleichskette
{\vergleichskette
{ K^2
}
{ = }{ K \times K
}
{  }{
}
{    }{
}
{   }{
}
}
{}{}{,}
die auf der Kurve liegen, die durch die Gleichung

\mavergleichskettedisp
{\vergleichskette
{ X^2Y  + 2Y^3+3Y^2
}
{ =} { 0
}
{ } {
}
{ } {
}
{ } {
}
}
{}{}{}
gegeben ist. Wie viele Lösungen gibt es?

}
{} {}


\inputaufgabegibtloesung
{}
{
Finde eine Gerade

\mavergleichskette
{\vergleichskette
{G
}
{ \subseteq }{ {\mathbb A}^{2}_{ {\mathbb C} }
}
{  }{
}
{    }{
}
{   }{
}
}
{}{}{,}
die die Kurve

\mavergleichskettedisp
{\vergleichskette
{ C
}
{ =} { V(X^3+Y^3 +1)
}
{ \subseteq} { {\mathbb A}^{2}_{ {\mathbb C} }
}
{ } {
}
{ } {
}
}
{}{}{}
in genau einem Punkt schneidet.

}
{} {}


\inputaufgabegibtloesung
{}
{
Zeige, dass die Neilsche Parabel

\mavergleichskettedisp
{\vergleichskette
{ C
}
{ =} { V \left(  Y^2-X^3  \right)
}
{ \subseteq} { {\mathbb A}^{2}_{ {\mathbb C} }
}
{ } {
}
{ } {
}
}
{}{}{}
jede Gerade durch den Punkt

\mavergleichskette
{\vergleichskette
{P
}
{ = }{ (1,1)
}
{ \in }{ C
}
{    }{
}
{   }{
}
}
{}{}{}
in mindestens einem weiteren Punkt trifft.

}
{} {}


\inputaufgabegibtloesung
{}
{
Begründe analytisch, dass es einen reellen Schnittpunkt des Einheitskreises
\mathl{V(x^2+y^2-1)}{} mit der Neilschen Parabel
\mathl{V( y^2-x^3)}{} gibt und bestimme numerisch die reelle $x$-Koordinate eines solchen Schnittpunktes mit einem Fehler $\leq 0,1$.

}
{} {}


\inputaufgabe
{}
{
Betrachte Gleichungen der Form

\mathdisp {y^2=G(x) \text{ mit } G(x) = x^3+ax^2+bx+c} {  }

über $\R$. Skizziere die verschiedenen Lösungsmengen für die Koeffizienten

\mavergleichskette
{\vergleichskette
{ a,b,c
}
{ \in }{ \{1,-1,0\}
}
{  }{
}
{    }{
}
{   }{
}
}
{}{}{.}

}
{} {}
Die folgende Aussage folgt über
Lemma 1.3
auch aus
Satz 6.1.


\inputaufgabe
{}
{
Es sei $K$ ein
\definitionsverweis {Körper}{}{}
und es sei
\maabbeledisp {} { K } { K^2
} { t } { (P(t),Q(t))
} {,}
eine durch zwei Polynome

\mavergleichskette
{\vergleichskette
{ P(t),Q(t)
}
{ \in }{ K[t]
}
{  }{
}
{    }{
}
{   }{
}
}
{}{}{}
gegebene Abbildung. Es sei $B$ das
\definitionsverweis {Bild}{}{}
dieser Abbildung und es sei

\mavergleichskette
{\vergleichskette
{ G
}
{ \subseteq }{ K^2
}
{  }{
}
{    }{
}
{   }{
}
}
{}{}{}
eine Gerade. Zeige, dass

\mavergleichskette
{\vergleichskette
{ B
}
{ \subseteq }{ G
}
{  }{
}
{    }{
}
{   }{
}
}
{}{}{}
ist oder dass der Durchschnitt
\mathl{B \cap G}{} endlich ist.

}
{} {}


\inputaufgabe
{}
{
Multipliziere in
\mathl{\Z[x,y,z]}{} die beiden Polynome

\mathdisp {x^5+3x^2y^2-xyz^3 \text{ und  } 2x^3yz+z^2+5xy^2z-x^2y} { . }


}
{} {}


\inputaufgabe
{}
{
Multipliziere in
\mathl{\Z/( 5 ) [x,y]}{} die beiden Polynome

\mathdisp {x^4+2x^2y^2-xy^3+2y^3 \text{ und  }  x^4y+4x^2y+3xy^2-x^2y^2+2y^2} { . }


}
{} {}


\inputaufgabe
{}
{
Es sei $R$ ein
\definitionsverweis {Integritätsbereich}{}{.}
Zeige, dass dann auch der
\definitionsverweis {Polynomring}{}{}
$R[X]$ integer ist.

}
{} {}


\inputaufgabegibtloesung
{}
{
Es sei $R$ ein
\definitionsverweis {Integritätsbereich}{}{}
und $R[X]$ der
\definitionsverweis {Polynomring}{}{}
über $R$. Zeige, dass die
\definitionsverweis {Einheiten}{}{}
von $R[X]$ genau die Einheiten von $R$ sind.

}
{} {}


\inputaufgabegibtloesung
{}
{
Es sei $K$ ein
\definitionsverweis {Körper}{}{.}
Zeige, dass die beiden folgenden Eigenschaften äquivalent sind:
\aufzaehlungzwei { $K$ ist
\definitionsverweis {algebraisch abgeschlossen}{}{.}
} {Jedes nicht-konstante Polynom

\mavergleichskette
{\vergleichskette
{ F
}
{ \in }{ K[X]
}
{  }{
}
{    }{
}
{   }{
}
}
{}{}{}
zerfällt in Linearfaktoren.
}

}
{} {}


\inputaufgabe
{}
{
Es sei $K$ ein
\definitionsverweis {algebraisch abgeschlossener Körper}{}{.}
Bestimme in
\mathl{K[X]}{} die
\definitionsverweis {irreduziblen Polynome}{}{.}

}
{} {}


\inputaufgabe
{}
{
Es sei $K$ ein
\definitionsverweis {algebraisch abgeschlossener Körper}{}{.}
Zeige, dass $K$ nicht endlich sein kann.

}
{} {}


  \zwischenueberschrift{Aufgaben zum Abgeben}


\inputaufgabe
{2}
{
Führe in
\mathl{\Z/( 7 ) [X]}{} folgende Polynomdivision aus.

\mathdisp {X^4+5X^2+ 3 \,  \text{ durch } \,  2X^2+X+6} { . }


}
{} {}


\inputaufgabe
{5}
{
Bestimme im \definitionsverweis {Polynomring}{}{}
\mathl{{\mathbb F}_{  3  } [X]}{}  alle normierten \definitionsverweis {irreduziblen Polynome}{}{} vom \definitionsverweis {Grad}{}{} $4$.

}
{} {}


\inputaufgabe
{3}
{
Bestimme alle Lösungen der Kreisgleichung

\mavergleichskettedisp
{\vergleichskette
{ x^2+y^2
}
{ =} { 1
}
{ } {
}
{ } {
}
{ } {
}
}
{}{}{}
für die Körper

\mavergleichskette
{\vergleichskette
{ K
}
{ = }{ \Z/( 2 )
}
{  }{
}
{    }{
}
{   }{
}
}
{}{}{,}

\mathl{\Z/( 5 )}{} und
\mathl{\Z/( 11)}{.}

}
{} {}


\inputaufgabe
{5}
{
Es sei

\mavergleichskette
{\vergleichskette
{ C
}
{ \subseteq }{ {\mathbb C}^2
}
{  }{
}
{    }{
}
{   }{
}
}
{}{}{}
das
\definitionsverweis {Bild}{}{}
unter der polynomialen Abbildung
\maabbeledisp {} { {\mathbb C} } { {\mathbb C}^2
} { t } { \left( t^3-t^2+4t+3  , \,   -t^2+5t-1                   \right)
} {.}
Bestimme ein Polynom

\mavergleichskette
{\vergleichskette
{ F
}
{ \neq }{ 0
}
{  }{
}
{    }{
}
{   }{
}
}
{}{}{}
in zwei Variablen derart, dass $C$ auf dem Nullstellengebilde zu $F$ liegt.

}
{} {}


\inputaufgabe
{4}
{
Wir betrachten die Abbildung
\maabbdisp {f} {\R} {S^1 \subseteq \R^2
} {,}
die einem Punkt

\mavergleichskette
{\vergleichskette
{ t
}
{ \in }{ \R
}
{  }{
}
{    }{
}
{   }{
}
}
{}{}{}
den eindeutigen Schnittpunkt
\mathl{\neq (0,-1)}{} der durch die beiden Punkte
\mathkor {} {(t,1)} {und} {(0,-1)} {}
gegebenen Geraden
\mathl{G_t}{} mit dem
\definitionsverweis {Einheitskreis}{}{}

\mavergleichskettedisp
{\vergleichskette
{ S^1
}
{ =} { \left\{ (x,y) \in \R^2  \mid   x^2+y^2 = 1        \right\}
}
{ } {
}
{ } {
}
{ } {
}
}
{}{}{}
zuordnet. Zeige, dass diese Abbildung wohldefiniert ist und bestimme die funktionalen Ausdrücke, die diese Abbildung beschreiben. Zeige, dass $f$
\definitionsverweis {differenzierbar}{}{}
ist. Ist $f$
\definitionsverweis {injektiv}{}{,}
ist $f$
\definitionsverweis {surjektiv}{}{?}

}
{} {}
